%Con handout eliminamos todos los \pause
\documentclass[handout, 14pt]{beamer}
%Paquetes.
\usepackage{graphicx}
\usepackage[spanish, mexico]{babel}
\usepackage{pgfpages}

%Etiquetas personalizables.
\newcommand{\tipo}[1]{\textcolor{blue}}

%\pgfpagesuselayout{4 on 1}
%Datos del documento.
\title{Curso de Arduino}
\author{Eliuth Montiel Navarrete}
\date{Septiembre 26 de 2026}

\usetheme{Madrid}

\begin{document}
\maketitle

\begin{frame}{Índice}
    \tableofcontents
\end{frame}

\begin{frame}
    
    \includegraphics[width=0.25\textwidth]{img/Eda.jpg}
    \frametitle{¿Quíén soy?}
    \textbf{Estudios}:
    Ingeniería mecatrónica, técnico en computación.
    \\
    \textbf{Cursos y diplomados}:
    Inteligencia artificial aplicada, diseño y programación de videojuegos, programación.
    \\
    \textbf{Gustos e intereses}:
    Videojuegos, programación, electrónica, idiomas, software libre, ciencia.

\end{frame}

\section{Clase 1: Fundamentos de programación}
\subsection{Variables}
\begin{frame}[fragile]{Variables}
    En C, una variable es un espacio de memoria en el cuál guardan datos para ser utilizados después. \cite{gfg_variables}
    \\
    Las variables se definen siguiendo la siguiente sintáxix:
    \begin{verbatim}
        tipo nombre = valor;
    \end{verbatim}
    \begin{exampleblock}{Ejemplos}
        \begin{verbatim}
        int edad = 45;
        float decimales = 1.15;
        char letra = 'a';
        bool verdad = true;
        \end{verbatim}
    \end{exampleblock}
\end{frame}

\begin{frame}[fragile]{Variables}
    Reglas para nombrar variables:
    \begin{itemize}
        \item Solo pueden contener letras, dígitos y guión bajo.
        \item No pueden comenzar con un dígito.
        \item No pueden ser \textbf{palabras reservadas}.
        \item El nombre debe ser único en el programa.
    \end{itemize}
\end{frame}

\subsection{Tipos de datos}
\begin{frame}{Tipos de datos}
    \begin{columns}{}
        \begin{column}{0.3\textwidth}
            \begin{block}{Básicos}
                \begin{itemize}
                    \item int
                    \item float
                    \item double
                    \item char
                    \item bool
                    \item void
                \end{itemize}
            \end{block}
        \end{column}

        \begin{column}{0.3\textwidth}
            \begin{block}{Derivados}
                \begin{itemize}
                    \item Arreglo
                    \item Apuntador
                \end{itemize}
            \end{block}
        \end{column}
        \begin{column}{0.3\textwidth}
            \begin{block}{Definidos por el usuario}
                \begin{itemize}
                    \item union
                    \item structure
                    \item enum
                \end{itemize}
            \end{block}
        \end{column}
    \end{columns}
    \cite{gfg_tipos}
\end{frame}

\begin{frame}[fragile]{Tipos de datos}
    \begin{exampleblock}{Ejemplos}
        \begin{verbatim}
            int i = 6;
            float f = 0.1;
            double d = 0.0005;
            char c = 'a';
            bool b = false;
            void v;
        \end{verbatim}
    \end{exampleblock}
\end{frame}

\begin{frame}[fragile]{Tipos de datos}
    \begin{exampleblock}{Ejemplos}
        \begin{verbatim}
    char nombre[5] = {'A','n','a','h','i'};
            char* charPointer = &nombre; 
        \end{verbatim}
    \end{exampleblock}
\end{frame}

\begin{frame}[fragile]{Tipos de datos}
    \begin{exampleblock}{Ejemplos}
        \begin{verbatim}
            union Dato{
                int x;
                float f;
            };
        \end{verbatim}
    \end{exampleblock}
\end{frame}

\begin{frame}[fragile]{Tipos de datos}
    \begin{exampleblock}{Ejemplos}
        \begin{verbatim}
            structure Estructura{
                char nombre[5];
                int edad;
                float calificacion;
            };
        \end{verbatim}
    \end{exampleblock}
\end{frame}

\begin{frame}[fragile]{Tipos de datos}
    \begin{exampleblock}{Ejemplos}
        \begin{verbatim}
enum Dia{
    LUNES, MARTE, MIERCOLES, JUEVES, VIERNES
};
        \end{verbatim}
    \end{exampleblock}
\end{frame}

\subsection{Asignación y expreciones}
\begin{frame}[fragile]{Asignación}
    Es una operación en la cual el valor del operando derecho
    se asigna al operando izquierdo (variable).\cite{microsoft_asignacion}
    \begin{exampleblock}{Ejemplos}
        \begin{verbatim}
                =, +=, -=
                *=, /=, %=
            <<=, >>=, &=, ^=, |=
        \end{verbatim}
    \end{exampleblock}
\end{frame}

\begin{frame}[fragile]{Expresiones}
    Es una secuencia de operadores y sus operandos, que especifica un cómputo.\cite{cpp_expressions}
    Se pueden clasificar como unarias, binarias, y otras.
    \begin{block}{Unarias}
        \begin{verbatim}
                ++. --, !
        \end{verbatim}
    \end{block}
    \begin{block}{Binarias}
        \begin{verbatim}
        +. -. ==, !=, <, > <=, >=
        \end{verbatim}
    \end{block}
\end{frame}

\section{Estructuras de control}
\subsection{Condicionales}
\begin{frame}{if}

\end{frame}

\begin{frame}{if-else}
    
\end{frame}

\begin{frame}{switch}
    
\end{frame}

\section{Desarrollo}
\begin{frame}{Mi segunda diapositiva}
    Esta es una lista de cosas por hacer:
    \begin{itemize}
        %\pause: Crea varias páginas con el contenido que está hasta \pause.
        \item Primera cosa. \pause
        \item Segunda cosa. \pause
        \item Tercera cosa. \pause
    \end{itemize}    

    Esta es una ecuación
    \begin{equation}
        e^{\pi i} + 1 = 0
    \end{equation}

\end{frame}

\section{Imágenes}
\begin{frame}{Imágenes}
    \begin{columns}
        \column{0.7\textwidth}
            \begin{figure}
                \centering
                \includegraphics{img/Eda.jpg}
                \caption{Eda pequeña}
                \label{fig:Eda-pequeña}
            \end{figure}
        \column{0.3\textwidth}
        Texto ejemplificando la superioridad de Eda pequeña.
    \end{columns}
\end{frame}

\begin{frame}{Bloques}
    \begin{block}{Texto importante}
        Esto es un texto importante.
    \end{block}
    \begin{alertblock}{Texto de alerta}
        Hola.\footnote{Esta es una nota.}
    \end{alertblock}
    \begin{exampleblock}{Texto de ejemplo}
        Esto es un ejemplo.
    \end{exampleblock}
\end{frame}

\section{Conclusión}
\begin{frame}
    Mis conclusiones son las siguientes:\\
    Sí.
\end{frame}

\section{Referencias}
\begin{frame}{Referencias}
    \bibliographystyle{IEEEtran}
    \bibliography{referencias.bib}
\end{frame}
\end{document}